\chapter*{Conclusiones}
\addcontentsline{toc}{chapter}{Conclusiones}

Las variedades algebraicas tienen una estructura bastante rígida pero, al mismo tiempo, bien comprendida en la geometría algebraica clásica. En contraste, las variedades complejas, aunque son más restrictivas que las variedades diferenciables reales, admiten una amplia gama de comportamientos globales. Por ejemplo, existen superficies complejas que no poseen funciones meromorfas no constantes (véase el ejemplo \ref{hopfsuperficie}), así como variedades complejas que no admiten métricas de Kähler (véase el contraejemplo de Hironaka \ref{hironakaejemplo}). Estos fenómenos nos muestran que las variedades complejas son mucho más diversas que las algebraicas, y en consecuencia, determinar si una variedad compleja es algebrizable y describir explícitamente las ecuaciones algebraicas que la definen, es una tarea lejos de ser trivial.

La algebrización de variedades complejas es un tema bastante profundo que admite múltiples enfoques. Por un lado, el teorema fundamental de Moishezon \ref{fundamentaldemoishezon} establece que una variedad compleja compacta es algebraica si, y solo si, es de Moishezon y de Kähler. Este resultado muestra claramente que la algebrización requiere al mismo tiempo de riqueza algebraica, medida mediante con la dimensión algebraica, y una riqueza geométrica, expresada mediante la existencia de una métrica de Kähler.

%Por otro lado, en este trabajo abordamos el problema de algebraización construyendo directa y explícitamente un encaje holomorfo de una variedad compleja compacta en un espacio proyectivo. Con esto conseguimos llevar toda la variedad, por medio de un biholomorfismo, al espacio proyectivo sin aplastarla o deformarla, es decir, preservando fielmente su estructura compleja. Adicionalmente, con el teorema de la aplicación propia (falta cita), esta incrustación termina siendo una subvariedad analítica del espacio proyectivo, concluyendo así por el teorema de Chow (falta cita) que nuestra variedad compleja es de hecho algebraica.

Por otro lado, en este trabajo abordamos el problema de algebraización construyendo directa y explícitamente un encaje holomorfo de una variedad compleja compacta en un espacio proyectivo. Con esto conseguimos llevar toda la variedad, por medio de un biholomorfismo, al espacio proyectivo sin aplastarla o deformarla, es decir, preservando fielmente su estructura compleja. Adicionalmente, con el teorema de la aplicación propia \ref{aplicacionpropia}, este encaje determina termina siendo una subvariedad analítica del espacio proyectivo, lo que nos conduce al teorema de Chow. Este resultado nos afirma que los subconjuntos analíticos del espacio proyectivo son conjuntos algebraicos, es decir, se pueden describir como conjuntos de ceros de una cantidad finita de polinomios homogéneos. De esta manera, concluimos que toda variedad compleja que admite un encaje holomorfo en un espacio proyectivo es algebrizable.

Además, en nuestra búsqueda de los encajes al proyectivo, encontramos que los haces lineales juegan un papel fundamental, ya que inducen una aplicación holomorfa a un espacio proyectivo con una base del espacio de secciones globales. Los haces lineales que inducen un encaje son llamados haces amplios, y así, concluimos que una variedad compleja es algebraica si, y solo si, admite un haz amplio (véase teorema \ref{proyectivossiamplio}). Como consecuencias, obtuvimos que toda superficie de Riemann es algebraica. Asimismo, ejemplificamos cómo la teoría de curvas elípticas permite describir explícitamente el toro complejo como una curva algebraica proyectiva. Finalmente, construimos explícitamente una realización algebraica de la variedad compleja $\mathbb P^1\times\mathbb P^1$, junto con la ecuación proyectiva que la describe.

Durante todo el trabajo observamos que el estudio de variedades complejas no puede abordarse de manera aislada, requiere interactuar con diversas áreas de las matemáticas. Por ejemplo, fue esencial comprender las propiedades básicas de las variedades algebraicas, estudiar la geometría local compleja y considerar el contexto clásico de las superficies de Riemann, cuya teoría en dimensión uno ilustra la profundidad y riqueza de la teoría general.

En esencia, la algebrización consiste en demostrar que la estructura analítica de una variedad compleja es suficientemente rica para reconstruirse a partir de objetos algebraicos. Esta tesis tiene como objetivo explicar y detallar el proceso de algebrización de variedades complejas. Para ello, se desarrolla un método de construir explícitamente encajes holomorfos al espacio proyectivo a través de haces lineales, y se detalla por qué esto es suficiente para que una variedad compleja admita una estructura algebraica. Esto motiva naturalmente la búsqueda y el estudio de condiciones que garanticen la existencia de haces amplios. Además, vimos la algebrización conecta directamente con resultados más profundos como los teoremas de Moishezon. Esta interacción entre geometría compleja y geometría algebraica constituye una de las ideas centrales desarrolladas a lo largo de este trabajo. Por lo tanto, el trabajo puede servir de base para futuras investigaciones sobre la interacción entre geometría compleja y geometría algebraica, ofreciendo una base sólida para explorar tópicos más avanzados en esta dirección.

%	En general, una variedad algebraica no está determinada necesariamente como puntos en un espacio proyectivo, por ejemplo una variedad $X$ definida como un esquema reducido separable de tipo finito sobre $\cc$. En esta caso, podemos obtener el espacio de puntos complejos $X(\cc)$, los cuales tienen la estructura topológica necesaria para obtener una variedad compleja \parencite{shafarevich2016basic}. 
%
%\begin{definition}
%	Sea $X$ una variedad algebraica sobre $\cc$. La \textbf{analiticación} de $X$, denotada por $X^{an}$, es la variedad compleja obtenida al considerar el conjunto de puntos complejos $X(\mathbb{C})$ dotado de la topología compleja y su estructura natural de variedad compleja.
%\end{definition}
%
%En general, se define el funtor
%\begin{equation*}
%	(-)^{an} : \mathrm{Var}*{\mathbb{C}} \longrightarrow \mathrm{Var}*{\mathbb{C}}^{an},
%\end{equation*}
%de la categoría de variedades algebraicas sobre $\cc$ a la categoría de variedades complejas\footnote{En general, a la categoría de espacios complejos \parencite{shafarevich2016basic}.}.